% answers_short.tex —— 2009 年·全国大纲卷Ⅱ·理科 · 简短答案速查
% 由 tools/make_short_answers.py 自 answers.tex 的 \gksolhead/\gkanswer 提取，
% 并与 2009/2009_数学真题_全国2理_答案要点.md 逐题比对后生成，勿手改（改源 MD 后重跑）。

\begin{gknotice}
\textbf{使用说明}：本节只给最终答案便于快速核对，完整推导见解析版。答案为编写者独立推导，\textbf{非官方标准答案}。
\end{gknotice}

\gkpart{选择题}
\begin{center}\begin{tabular}{ccccccccccccc}
\toprule
题号 & 1 & 2 & 3 & 4 & 5 & 6 & 7 & 8 & 9 & 10 & 11 & 12 \\
\midrule
答案 & A & B & C & B & C & C & C & D & D & C & A & A \\
\bottomrule
\end{tabular}\end{center}

\gkpart{填空题}
\begin{itemize}[leftmargin=2.6em]
  \item[\textbf{13.}] $6$
  \item[\textbf{14.}] $9$
  \item[\textbf{15.}] $8$
  \item[\textbf{16.}] $5$
\end{itemize}

\gkpart{解答题}
\begin{itemize}[leftmargin=2.6em]
  \item[\textbf{17.}] $B=60^\circ$。
  \item[\textbf{18.}] (1) 见解析版；(2) $B_1C$ 与平面 $BCD$ 所成的角为 $30^\circ$
  \item[\textbf{19.}] (1) $b_n=3\cdot2^{n-1}$；(2) $a_n=(3n-1)2^{n-2}$
  \item[\textbf{20.}] (1) 甲组抽取 $2$ 人，乙组抽取 $1$ 人；(2) $\dfrac8{15}$；(3) $\Pr(\xi=0)=\dfrac15$，$\Pr(\xi=1)=\dfrac{34}{75}$，$\Pr(\xi=2)=\dfrac{22}{75}$，$\Pr(\xi=3)=\dfrac4{75}$，$\mathbb E\xi=\dfrac65$
  \item[\textbf{21.}] (1) $a=\sqrt3,\ b=\sqrt2$；(2) 存在 $P\left(\dfrac32,\dfrac{\sqrt2}{2}\right)$、$P\left(\dfrac32,-\dfrac{\sqrt2}{2}\right)$；相应直线为 $y=\sqrt2(x-1)$、$y=-\sqrt2(x-1)$
  \item[\textbf{22.}] (1) $0<a<\dfrac12$；在 $\left(-1,\dfrac{-1-\sqrt{1-2a}}2\right)$ 和 $\left(\dfrac{-1+\sqrt{1-2a}}2,+\infty\right)$ 上单调递增，在两根之间单调递减；(2) $f(x_2)>\dfrac{1-2\ln2}{4}$
\end{itemize}
